\section{Background and related work}
\label{sec:background}

\paragraph{Weight-only quantization at batch 1.} When a model serves one
request at a time, each new token multiplies every weight matrix by a single
vector. So the bytes of weights the GPU reads
set the time per token. GPTQ~\citep{gptq2023} quantizes a matrix one column at
a time. It pushes each rounding error onto the columns not yet quantized, using
second-order statistics from calibration text. AWQ~\citep{awq2024} rescales
input channels by activation size before rounding to 4 bits. With groups of 128
weights it is the usual 4-bit baseline. vLLM serves it with the Marlin
kernel~\citep{vllm2023,marlin2024}.

\paragraph{Two-bit vector quantization.} Below three bits, rounding each weight
on its own loses too much. The strong methods encode several weights at once,
as a choice from a codebook. QuIP\#~\citep{quipsharp2024} uses the $E_8$
lattice. QTIP~\citep{qtip2024} uses a trellis. Its hybrid code hashes each
state to index a 2~KiB table, and its computed codes, such as 3INST, need no
table. AQLM and VPTQ~\citep{aqlm2024,vptq2024} learn their codebooks, of at
most $2^{16}$ entries. Their kernels decode by reading these tables, and both
papers tie the kernel's speed to how well the tables fit in the GPU cache.
QuIP\# and QTIP first apply a random Hadamard transform~\citep{quipsharp2024},
a rotation that spreads outliers across
coordinates~\citep{quip2023,quarot2024}. llama.cpp~\citep{llamacpp} runs
models on local machines. Its IQ2 formats decode from grids of 256 to
1{,}024 points held in a lookup table.

\paragraph{The Leech lattice.} \citet{llvq2026} quantize the weights in blocks
of 24 on $\Lambda_{24}$, the densest sphere packing in 24
dimensions~\citep{cohn2017}. They find that their codebook makes the model
less sensitive to rotating the weights first. Inside a GPTQ-style loop, each block is coded as a
lattice direction (the shape) and a length (the gain). A block takes 48 bits,
or 2 bits per weight. In their best codebook on Gaussian data, 47 bits pick one
of $1.1\times10^{14}$ lattice points inside a ball, and one bit gives the
gain. The model results we cite use a larger ball: all 48 bits pick one of
$2.8\times10^{14}$ points, and no bit goes to the gain. With it, on Qwen3-4B
they report 60.7 on MMLU without fine-tuning and 62.8 with it. In the same two
settings they report 57.4 and 59.5 for QTIP with its 3INST code, and 48.6 and
52.9 for QuIP\#. Their GPU kernel,
however, decodes only one shell of the lattice (the points at one distance from
the origin). Those scores come from a codebook with several shells.

\paragraph{Our earlier kernel.} In \citet{llvq1preprint} we served their
47-bit codebook with its gain bit. At load time, each 47-bit index and its gain
bit were unfolded into a 112-bit record, the \Planes{} layout. The record holds a class, a gain bit, a sign mask and
three bit planes. The class names a set of magnitudes, and the planes pick one
for each weight. A fused kernel reads these records and multiplies in one pass.
It reads 4.804 bits per weight. This paper replaces that layout.

\paragraph{Trellises of the Golay code.} $\Lambda_{24}$ is built from the
binary Golay code $\mathcal{G}_{24}$. When its coordinates are cut into three
octads, codewords with eight ones that share no position, this code has a
trellis with 64 states at each cut~\citep{forney1988}. Such trellises lead to
efficient maximum-likelihood decoding: finding the codeword or lattice point
most likely to have been sent~\citep{forney1988}. \citet{vardy1991} decode the
Golay code by maximum likelihood through another structure, the hexacode. We
use the trellis to give every entry of a quantizer codebook an address.
